\chapter{Konzept}
\label{cha:Konzept}

In diesem Kapitel wird ein Umsetzungskonzept für die im letzten Kapitel beschriebenen Anforderungen erarbeitet.


Umsetzung dll / lib / oder exe?

\section{Struktur der Programmbibliothek}
\subsection{Verwalten von Isolates und Contexts}
\subsection{Freilegen des DebugInspector Protokolls}
\subsection{Dynamische Parametrierung}

\section{Struktur des V8-API-Agenten}
\subsection{Parametrierung}

\section{Umsetzung}
\subsection{Zu verwendende Compiler}
\subsection{Benötigte Abhängigkeiten}
\subsection{Build Strategie}

\subsection{Compiler}

wahl der compiler auf linux und windows

\subsection{Handbuch Integratoren}

Typedoc / Doxygen.

\subsection{Linken externer Bibliotheken}

Im folgenden Abschnitt wird erklärt, wie ein Programm mit externen Bibliotheken umgehen kann. Die verschiedenen Möglichkeiten
werden dann abgewogen und in einem Fazit wird entschieden, welche Methode für den Prototyp verwendet wird.

\subsubsection{Statisches Linken}

Wird eine Bibliothek statisch gelinkt, dann werden die Funktionen dieser Bibliothek direkt in die ausführbare Datei kopiert.
Voraussetzung dafür ist, dass die gelinkte Bibliothek selbst statisch ist. Dies bedeutet, dass mindestens eine Objektdatei mithilfe des
Kommandozeilenprogrammes \textit{ar (archiver utility)} in ein Archiv umgewandelt wurde. \textit{Ar} fügt zu einem solchen Archiv eine Tabelle hinzu, in der vermerkt ist welche Symbole, wo im Archiv zu finden sind \cite{staticLinking}.

\subsubsection{Dynamisches Linken}

Beim dynamischen Linken werden die Funktionen einer dynamische Bibliothek in der ausführbaren Datei nur referenziert. Die dynamische
Bibliothek wird erst dann geladen, wenn das Programm gestartet wird. Ein wesentlicher Vorteil von dynamischen Bibliotheken gegenüber statischen
ist, dass der ausführbare Code für alle Prozesse, die diese Bibliothek benutzen nur einmal in den Speicher geladen wird \cite{dynLinking}.

\subsubsection{Fazit}

Da der Prototyp nur als einzelner Prozess laufen soll, ist durch die Verwendung von dynamischen Bibliotheken kein wirklicher Vorteil gegeben.
Es macht daher mehr Sinn alle externen Bibliotheken statisch in die ausführbare Datei zu linken, da man somit von externen Bibliotheken
unabhängig bleibt.
